\begin{abstract}
Antimicrobial peptides (AMPs) are promising candidates for next-generation
antibiotics, but computational AMP discovery remains limited by saturated
binary benchmarks, assay-dependent potency measurements, and conditional
generators whose controllability is rarely audited property by property.
We study these limitations with \textbf{JEPA-AMP}, a compact
Joint-Embedding Predictive Architecture pre-trained on \textbf{868{,}724}
AMP sequences and evaluated as a shared backbone for classification,
MIC and HC50 regression, and three conditional generation paradigms.

Across AMPlify, GRAMPA, and QMAP, AMP-specific pre-training gives useful but
limited gains.
JEPA-AMP approaches strong supervised baselines on AMP classification
(AUROC~\textbf{0.958}, MCC~\textbf{0.802}) without surpassing the published
AMPlify ensemble, and improves GRAMPA MIC regression over an ESM-2~(35M)
baseline (best JEPA Pearson~$r{=}\mathbf{0.640}$, RMSE~\textbf{0.619}).
On the homology-aware QMAP benchmark, JEPA-AMP is competitive with the
best reported full E.~coli result (PCC~\textbf{0.512}) and distinctly stronger
on high-efficiency E.~coli (PCC~\textbf{0.388} vs.\ prior best 0.290, $+34\%$).

A frozen embedding comparison reveals a \emph{plasticity} gap.
ESM-2 frozen embeddings score higher on $k$-NN MIC prediction and
AMP linear-probe AUROC, yet JEPA-AMP performs better after supervised
fine-tuning. This suggests that JEPA pre-training keeps useful capacity
that appears after task-specific tuning.

For generation, the proposed dual-pathway decoder---combining per-layer AdaLN
conditioning with a cross-attention condition token---achieves reliable charge
control ($R^2{=}\mathbf{0.866}$ vs.\ negative $R^2$ for weaker baselines),
while GRAVY, length, and species-selective MIC control remain limited.
We also compare three generator settings on the same frozen JEPA backbone:
the AR baseline, a non-autoregressive (NAR) decoder, and a masked diffusion
decoder. The AR decoder gives the best charge control, while masked diffusion
is the strongest new paradigm (charge $R^2{=}0.702$, length $R^2{=}1.000$).
An extended 7-condition AR model gives partial extra control for helix
($R^2{=}0.533$) and pI ($R^2{=}0.330$), but weak control for other added targets.
MIC-conditioned generation shifts model-predicted global activity relative to
matched physicochemical controls (broad-spectrum:
$\Delta\log_2\text{MIC}{=}{-}0.250$ under JEPA oracle, $-0.291$ under
independent ESM-2 scorer), but these are \emph{oracle-mediated dry-lab screens},
not biological validation.
All code, configs, and locked artifacts are released at
\url{https://github.com/[repo]}.
\end{abstract}
